% !TeX root = ../main.tex
% Add the above to each chapter to make compiling the PDF easier in some editors.

\chapter{Theoretical Background}\label{chapter:theoretical_background}

\section{Databases and Query Execution}

A \emph{database management system} (DBMS) is software that stores, organizes, and retrieves structured data.
Users interact with such systems primarily through \emph{Structured Query Language} (SQL), 
a declarative language that lets users specify what data they want rather than how to retrieve it~\parencite{silberschatz2019}.
For example, a simple SQL query might select all orders placed by a particular customer by joining an \texttt{orders} table with a \texttt{customers} table and filtering by a customer name.
Because SQL describes the desired result and not the retrieval procedure, the DBMS is free to choose any execution strategy that produces the correct answer.

When a DBMS receives an SQL query, it processes the query in three successive phases~\parencite{silberschatz2019}:

\begin{enumerate}
  \item \textbf{Parsing and translation.}
  The SQL text is checked for syntactic correctness and translated into an internal relational-algebra expression.
  During this step, the parser also expands any \emph{view} references — a view is a named, stored query that acts as a virtual table — by substituting them with their underlying relational-algebra definition.

  \item \textbf{Optimization.}
  The \emph{optimizer} takes the relational-algebra expression and uses statistical summaries of the stored data to evaluate many possible \emph{execution plans}.
  A plan is a tree of physical \emph{operators} — such as table scans, joins, sorts, and aggregations — along with instructions for how each operator should be executed (e.g., whether to use an index).

  A key input to the optimizer is the \emph{cardinality} of each operator: the number of rows (tuples) it is expected to produce or consume.
  Since the query has not yet been executed, the optimizer relies on \emph{estimated cardinalities}, derived from table sizes and column value distributions.
  These estimates directly influence which plan the optimizer considers cheapest and therefore selects.

  The \emph{actual cardinality} — the true number of rows processed at runtime — only becomes known after execution.
  Estimated and actual cardinalities can differ substantially, which is a recurring challenge in performance prediction~\parencite{leis2018job}.

  \item \textbf{Evaluation.}
  The query-execution engine runs the chosen execution plan against the stored data and produces the query result.
\end{enumerate}

The query plan is central to performance, since the optimizer's choice of operators and their ordering directly determines how much work the DBMS has to perform.
Different DBMSs implement these phases differently, and especially the execution step varies widely.
\emph{PostgreSQL}, for instance, executes plans using an iterator-based model described below.
\emph{Umbra}, by contrast, compiles query plans into native machine code at runtime.
These architectural choices have a large impact on query execution time~\parencite{neumann2011}.

\subsection{Iterator-Based Execution (PostgreSQL)}

PostgreSQL (PG) uses the \emph{iterator model} for query execution, which is also known as the \emph{Volcano model} or \emph{pull-based execution} after the system that popularized it~\parencite{graefe1994}.

In this model, every operator in the query plan tree exposes a uniform three-call interface: \texttt{open()}, \texttt{next()}, and \texttt{close()}.
The central call is \texttt{next()}: when invoked, the operator produces and returns the next output tuple (i.e., one row of data).
If no more tuples are available, the operator signals end-of-stream.
The \texttt{open()} call initializes the operator and its state, and \texttt{close()} releases any resources held.

Operators are composed into a tree that mirrors the query plan.
The root operator drives execution: each time it needs another tuple, it calls \texttt{next()} on its child operator.
That child, in turn, calls \texttt{next()} on its own children, and so on, all the way down to the leaf operators — typically table scans — which actually read data from disk or memory.
Tuples thus \emph{flow upward} through the tree one at a time, on demand.
This is the defining property of pull-based execution: no operator pushes data forward; instead, every operator waits to be asked.

As a concrete example, consider a query that scans a table, filters rows by a predicate, and then sorts the result.
The tree would consist of a sort operator at the root, a filter operator as its child, and a sequential scan at the leaf.
When execution starts, the sort operator repeatedly calls \texttt{next()} on the filter, which in turn calls \texttt{next()} on the scan, until all matching rows have been collected.
Only then does the sort operator produce its first output tuple.

PG exposes this plan tree through the \texttt{EXPLAIN} and \texttt{EXPLAIN ANALYZE} commands.
They are prepended to any SQL statement and cause PG to return the query plan instead of, or in addition to, the query result.
\texttt{EXPLAIN} alone shows the chosen plan with the optimizer's cost and cardinality estimates for each node, without executing the query.
\texttt{EXPLAIN ANALYZE} additionally executes the query and annotates every node with the actual number of rows produced and the actual time spent~\parencite{postgresql_docs_query_path}.
This makes \texttt{EXPLAIN ANALYZE} output the primary source of per-operator timing information in PG and is also the format used by the ZeroShot benchmark that this thesis builds on~\parencite{zeroshot}.

A key advantage of the iterator model is its simplicity and composability.
New operators can be added independently, as long as they implement the three-call interface~\parencite{graefe1994}.
However, this design also has a well-known downside: calling \texttt{next()} individually for every single tuple introduces per-tuple function call overhead and hurts CPU branch prediction and cache behavior~\parencite{neumann2011}.
This is one reason why more recent systems like Umbra take a different approach, as described in the next subsection.

\subsection{Compiled Query Execution (Umbra)}

Umbra is a relational DBMS developed at the Technical University of Munich that is designed for high performance on modern hardware~\parencite{umbra2020}.
Unlike PG's interpreter-based approach, Umbra takes a fundamentally different execution strategy: it compiles each query plan into native machine code at query time~\parencite{umbra2020, neumann2011}.

The motivation for this is to eliminate the per-tuple overhead of the iterator model.
Recall that in PG's pull-based execution, every single tuple causes a chain of \texttt{next()} calls to travel through the operator tree.
This introduces substantial function call overhead and makes it difficult for the CPU to predict branches and keep data in cache.
In a compiled system, these operator boundaries dissolve.
Instead of calling functions for each tuple, Umbra generates a single piece of native machine code that fuses the logic of multiple operators into tight loops.
Data is \emph{pushed} forward through the code rather than pulled upward through function calls.

Concretely, the compilation model works as follows~\parencite{neumann2011}.
Each operator participates in a \texttt{produce}/\texttt{consume} interface, which is a code-generation protocol invoked once at query compile time.
Code generation starts at the root, which calls \texttt{produce()} on its child; this propagates downward until the leaf scan is reached.
The scan emits a \texttt{for}-loop over its tuples and, for each tuple, calls \texttt{consume()} on its parent.
Each \texttt{consume()} invocation emits that operator's processing logic inline — a predicate check, a buffer write, etc.\ — and then calls \texttt{consume()} further up.
The result is a compact code fragment: a single \texttt{for}-loop from the scan, with the logic of every subsequent operator fused inside it, and no function-call boundaries between operators.

In the scan-filter-sort example from the previous subsection, the first fragment is exactly such a loop.
It scans tuples, checks the filter predicate inline, and writes qualifying tuples into a sort buffer.
The sort, however, cannot produce output until all input is accumulated.
Therefore, it acts as a materialization boundary, after which a second fragment handles the sorted output.
These materialization boundaries between fragments are the concept of \emph{pipeline breakers}, introduced in the next subsection.

At runtime, only the generated code runs.
The scan drives execution by iterating over tuples, and each tuple is processed through the entire pipeline inline.
This is the push-based contrast with PostgreSQL: in PG, the top-most operator drives execution by issuing \texttt{next()} calls all the way down the tree for every single tuple at runtime.
In Umbra, those operator boundaries are compiled away.
The generated code is handed to the LLVM compiler framework~\parencite{neumann2011}, which produces native machine code in typically a few milliseconds.
This overhead is amortized for non-trivial queries.

As a concrete consequence for this thesis: \emph{T3}, the performance prediction model this work builds on, was originally designed for Umbra and therefore relies on Umbra's query plans and their structure~\parencite{t3}.
The key structural concept that Umbra's compilation naturally exposes is the notion of \emph{pipelines}, which is described next.

\subsection{Pipelines and Pipeline Breakers}

The code fragments described above correspond to what are formally called \emph{pipelines} — a central concept in compiled query execution and the structural basis of T3's prediction approach~\parencite{t3}.

A \emph{pipeline} is a maximal sequence of operators that can process tuples in a fully streaming fashion, without any need to fully materialize an intermediate result.
Within a pipeline, each tuple produced by the scan at the bottom flows immediately through all subsequent operators — being filtered, transformed, or joined — without waiting for any other tuple.

Moreover, a \emph{pipeline breaker} is an operator that interrupts this streaming flow, because it must see \emph{all} of its input before it can produce any output.
The following are typical examples of pipeline breakers~\parencite{t3}:

\begin{itemize}
  \item \textbf{Sort.} A sort operator cannot emit a single sorted tuple until it has collected and sorted all input tuples.
  \item \textbf{Aggregation.} A group-by operator must accumulate all input tuples into groups before it can emit any aggregated result.
  \item \textbf{Hash join (build side).} The build side of a hash join must be fully read into a hash table before any tuple from the probe side can be matched.
    The build side therefore forms its own pipeline, terminating at the hash table.
\end{itemize}

Not every operator is a pipeline breaker.
A filter or a projection passes each tuple through without any materialization.
The probe side of a hash join — where tuples are looked up in an already-built hash table — behaves the same way.
Therefore, a single pipeline can contain multiple consecutive join probe stages.

Thus, a query plan decomposes into a directed acyclic graph (DAG) of pipelines.
Each pipeline starts either at a base table scan or by reading the materialized output of a prior pipeline breaker, processes tuples in a streaming fashion, and ends either at the next pipeline breaker or at the final output.


\section{Database Performance Prediction}

\subsection{Benchmarking}


sowas auch schreiben wie ssb, tpch, etc.

\subsection{Error Metrics}
q-error und percentile.

\subsection{Traditional Cost Models and Learned Approaches}

\subsection{Neural Networks and Decision Trees}

\subsection{T3: Tuple Time Tree}

This decomposition is fundamental to T3's performance prediction approach~\parencite{t3}.
Rather than predicting the runtime of a whole query in a single step, T3 predicts the runtime of each pipeline individually and sums the per-pipeline predictions to obtain the total query runtime.
Intuitively, each pipeline represents a distinct computational unit with its own characteristic workload — its scan cardinality, the operators it contains, and the fraction of tuples that reach each operator.
Predicting per-pipeline is therefore simpler and more accurate than predicting for the whole query at once.
The pipelines and their structure are also the basis for the feature vectors used by T3, as discussed in the Methodology chapter.

\section{Related Work}

\section{--------------------------------}

\section{Database Basics}

operator type, estimated/actual cardinalities, widths, and runtimes

\section{Performance Prediction}

draft: beschreiben, was performance prediction ist, was die ziele sind, etc.
benchmarking, etc.
Cost Model beschreiben

\section{Benchmarking}
timed queries, etc.


\section{Cost Models}

\section{Decision Trees}

\section{PostgreSQL}

\section{Umbra}

\section{T3}


\section{Section}
Citation test~\parencite{latex}.

Acronyms must be added in \texttt{main.tex} and are referenced using macros. The first occurrence is automatically replaced with the long version of the acronym, while all subsequent usages use the abbreviation.

E.g. \texttt{\textbackslash ac\{TUM\}, \textbackslash ac\{TUM\}} $\Rightarrow$ \ac{TUM}, \ac{TUM}

For more details, see the documentation of the \texttt{acronym} package\footnote{\url{https://ctan.org/pkg/acronym}}.
\subsection{Subsection}

See~\autoref{tab:sample}, \autoref{fig:sample-drawing}, \autoref{fig:sample-plot}, \autoref{fig:sample-listing}.

\begin{table}[htpb]
  \caption[Example table]{An example for a simple table.}\label{tab:sample}
  \centering
  \begin{tabular}{l l l l}
    \toprule
      A & B & C & D \\
    \midrule
      1 & 2 & 1 & 2 \\
      2 & 3 & 2 & 3 \\
    \bottomrule
  \end{tabular}
\end{table}

\begin{figure}[htpb]
  \centering
  % This should probably go into a file in figures/
  \begin{tikzpicture}[node distance=3cm]
    \node (R0) {$R_1$};
    \node (R1) [right of=R0] {$R_2$};
    \node (R2) [below of=R1] {$R_4$};
    \node (R3) [below of=R0] {$R_3$};
    \node (R4) [right of=R1] {$R_5$};

    \path[every node]
      (R0) edge (R1)
      (R0) edge (R3)
      (R3) edge (R2)
      (R2) edge (R1)
      (R1) edge (R4);
  \end{tikzpicture}
  \caption[Example drawing]{An example for a simple drawing.}\label{fig:sample-drawing}
\end{figure}

\begin{figure}[htpb]
  \centering

  \pgfplotstableset{col sep=&, row sep=\\}
  % This should probably go into a file in data/
  \pgfplotstableread{
    a & b    \\
    1 & 1000 \\
    2 & 1500 \\
    3 & 1600 \\
  }\exampleA
  \pgfplotstableread{
    a & b    \\
    1 & 1200 \\
    2 & 800 \\
    3 & 1400 \\
  }\exampleB
  % This should probably go into a file in figures/
  \begin{tikzpicture}
    \begin{axis}[
        ymin=0,
        legend style={legend pos=south east},
        grid,
        thick,
        ylabel=Y,
        xlabel=X
      ]
      \addplot table[x=a, y=b]{\exampleA};
      \addlegendentry{Example A}
      \addplot table[x=a, y=b]{\exampleB};
      \addlegendentry{Example B}
    \end{axis}
  \end{tikzpicture}
  \caption[Example plot]{An example for a simple plot.}\label{fig:sample-plot}
\end{figure}

\begin{figure}[htpb]
  \centering
  \begin{tabular}{c}
  \begin{lstlisting}[language=SQL]
    SELECT * FROM tbl WHERE tbl.str = "str"
  \end{lstlisting}
  \end{tabular}
  \caption[Example listing]{An example for a source code listing.}\label{fig:sample-listing}
\end{figure}
